\BourbakiExercisesSection

\begin{enumerate}
\def\labelenumi{\arabic{enumi}.}
\tightlist
\item
  Let E be an A-algebra with a basis of two elements \(e_1, e_2\) and a unit element \(e = \lambda e_1 + \mu e_2\) .
\end{enumerate}

\begin{enumerate}
\def\labelenumi{(\alph{enumi})}
\item
  Show that the ideal \(a\) of A generated by \(\lambda\) and \(\mu\) is the whole ring A (in the contrary case, note that \(\mathrm{E} / \mathrm{aE} = \mathrm{E} \otimes_{\mathbf{A}} (\mathrm{A} / \mathrm{a}) = \{0\}\) would hold contrary to the fact that \(\mathrm{A} / \mathrm{a} \neq \{0\}\) and E is a free A-module).
\item
  If \(\alpha, \beta\) are two elements of A such that \(\alpha\lambda + \beta\mu = 1\) , show that the elements \(e\) and \(u = \beta e_1 - \alpha e_2\) also form a basis of E and deduce that E is a quadratic A-algebra.
\end{enumerate}

\begin{enumerate}
\def\labelenumi{\arabic{enumi}.}
\setcounter{enumi}{1}
\item
  Show that every quadratic algebra over \(\mathbf{Z}\) is isomorphic to an algebra of type \((0, n)\) with \(n \in \mathbf{N}\) , or an algebra of type \((m, 0)\) with \(m\) a non-square, or an algebra of type \((m, 1)\) , where \(m\) is an integer which is not of the form \(k(k - 1)\) ; moreover, no two of these algebras are isomorphic.
\item
  \begin{enumerate}
  \def\labelenumii{(\alph{enumii})}
  \tightlist
  \item
    Let K be a commutative field. Show that, for K to be the centre of a quaternion algebra of type \((\alpha, \beta, \gamma)\) over K, it is necessary and sufficient that one of the following cases hold: (1) \(\beta\gamma \neq 0\) ; (2) \(\gamma = 0\) , \(\beta \neq 0\) and \(\beta^{2} + 4\alpha \neq 0\) ; (3) \(\beta = 0\) , \(\alpha \neq 0\) or \(\gamma \neq 0\) and \(2 \neq 0\) in K.
  \end{enumerate}
\end{enumerate}

\begin{enumerate}
\def\labelenumi{(\alph{enumi})}
\setcounter{enumi}{1}
\item
  Let K be a commutative field such that \(2 \neq 0\) in K; show that the quaternion algebra over K of type \((1, \gamma)\) is isomorphic to the matrix algebra \(\mathbf{M}_{2}(\mathbf{K})\) (consider the basis of this algebra consisting of the elements \(\frac{1}{2}(1 + i)\) , \(\frac{1}{2}(1 - i)\) , \((j + k/2\gamma)\) , \(\frac{1}{2}(j - k)\) ).
\item
  Let K be a commutative field such that \(2 \neq 0\) in K and let E be the quaternion algebra over K of type \((\alpha, \gamma)\) . A quaternion \(z \in E\) is called pure if \(\overline{z} = -z\) (or also \(T(z) = 0\) ). Show that if z is pure, so is every quaternion \(tzt^{-1}\) , where \(t \in E\) is invertible.
\item
  If \(u, v\) are any two quaternions, \(uv - vu\) is pure. Deduce that if \(u, v, w\) are any three quaternions, then
\end{enumerate}

\[
(u v - v u) ^ {2} w = w (u v - v u) ^ {2}.
\]

\begin{enumerate}
\def\labelenumi{(\alph{enumi})}
\setcounter{enumi}{4}
\tightlist
\item
  Show that if there exists in E a quaternion \(z \neq 0\) such that \(\mathrm{N}(z) = 0\) , there also exists a pure quaternion \(z' \neq 0\) such that \(\mathrm{N}(z') = 0\) . (Note, in the notation of formula (33) (no. 5), that, if \(\alpha\) is not a square in K, the existence of \(z \neq 0\) in E such that \(\mathrm{N}(z) = 0\) is equivalent to the existence of an element \(y \in \mathrm{K} + \mathrm{Ki}\) such that \(\gamma = \mathrm{N}(y)\) .)
\end{enumerate}

\begin{enumerate}
\def\labelenumi{\arabic{enumi}.}
\setcounter{enumi}{4}
\tightlist
\item
  Over a commutative field K in which 2 = 0, every quaternion algebra E of type \((\alpha, 0, \gamma)\) is commutative and the square of every \(x \in E\) belongs to K; the subalgebra \(\mathrm{K}(x)\) of E generated by an element \(x \in K\) is therefore a quadratic algebra of type \((\lambda, 0)\) over K and E is a quadratic algebra over \(\mathrm{K}(x)\) .
\end{enumerate}

Show that the set C of \(x \in E\) whose square is equal to the square of an element of K is a vector subspace of E whose dimension is equal to 1, 2 or 4. If C is of dimension 1 (in which case C = K), E is a field. If C is of dimension 2, there exists a quadratic algebra \(\mathrm{K}(x)\) contained in E which is a field and E has a basis over \(\mathrm{K}(x)\) consisting of two elements 1 and u with \(u^{2} = 0\) ; the set of \(y \in E\) such that \(y^{2} = 0\) is an ideal a of dimension 2 over K and E/a is isomorphic to \(\mathrm{K}(x)\) . Finally, if C is of dimension 4 (and hence equal to E), the set a of \(y \in E\) such that \(y^{2} = 0\) is an ideal of dimension 3 and E/a is isomorphic to K; there exists in E a basis \((1, e_{1}, e_{2}, e_{3})\) such that \(e_{1}^{2} = e_{2}^{2} = e_{3}^{2} = 0\) , \(e_{1}e_{2} = e_{3}\) , \(e_{1}e_{3} = e_{2}e_{3} = 0\) ; \(Ke_{3} = b\) is the only ideal of dimension 1 in E, b is the annihilator of a and a/b is the direct sum of two ideals of E/b, of dimension 1, which annihilate one another.

\begin{enumerate}
\def\labelenumi{\arabic{enumi}.}
\setcounter{enumi}{5}
\tightlist
\item
  Let \(\mathbf{K}\) be a commutative field in which \(2 \neq 0\) and \(\mathbf{E}\) a quaternion algebra over \(\mathbf{K}\) of type \((0, 0, \gamma)\) .
\end{enumerate}

\begin{enumerate}
\def\labelenumi{(\alph{enumi})}
\item
  If \(\gamma\) is not a square in \(\mathbf{K}\) , there exists no left (resp. right) ideal in \(\mathbf{E}\) of dimension 1 over \(\mathbf{K}\) ; the set \(a\) of \(x \in \mathbf{E}\) such that \(x^2 = 0\) is a two-sided ideal of dimension 2 over \(\mathbf{K}\) and \(\mathbf{E} / a\) is a field, a quadratic algebra over \(\mathbf{K}\) .
\item
  If \(\gamma\) is a square \(\neq 0\) in \(\mathbf{K}\) , there exists in \(\mathbf{E}\) a basis \((e_i)_{1\leq i\leq 4}\) with the following multiplication table:
\end{enumerate}

\(e_1\)

\(e_2\)

\(e_3\)

\(e_4\)

\(e_1\)

\(e_1\)

0

\(e_3\)

0

\(e_2\)

0

\(e_2\)

0

\(e_4\)

\(e_3\)

0

\(e_3\)

0

0

\(e_4\)

\(e_4\)

0

0

0

The set a of \(x \in E\) such that \(x^{2} = 0\) is a two-sided ideal of dimension 2, the direct sum of the two two-sided ideals \(Ke_{3}\) and \(Ke_{4}\) which annihilate one another; the latter are the only (left or right) ideals of dimension 1 in E. The quotient algebra E/a is isomorphic to the product of two fields isomorphic to K.

\begin{enumerate}
\def\labelenumi{(\alph{enumi})}
\setcounter{enumi}{2}
\tightlist
\item
  If \(\gamma = 0\) , the set \(a\) of \(x \in E\) such that \(x^2 = 0\) is a two-sided ideal of dimension 3; \(Kk = b\) is the only (left or right) ideal of dimension 1 in \(E\) ; it is a two-sided ideal, the left and right annihilator of \(a\) . In the quotient algebra \(E/b\) , \(a/b\) is the direct sum of two two-sided ideals of dimension 1 which annihilate one another; finally \(E/a\) is a field isomorphic to \(K\) .
\end{enumerate}

\begin{enumerate}
\def\labelenumi{\arabic{enumi}.}
\setcounter{enumi}{6}
\tightlist
\item
  Let K be a commutative field in which \(2 \neq 0\) and E the algebra over K with a basis of four elements 1, i, j, k (1 the unit element) with multiplication table
\end{enumerate}

\[
i ^ {2} = j ^ {2} = k ^ {2} = 1, \quad i j = j i = k, \quad j k = k j = i, \quad k i = i k = j.
\]

Show that E is isomorphic to the product of four fields isomorphic to K (consider the basis of E consisting of the elements \((1 + \varepsilon i)(1 + \varepsilon'j)\) , where \(\varepsilon\) and \(\varepsilon'\) are equal to 1 or -1).

The algebra E is the algebra (over K) of the product of two cyclic groups of order 2. Generalize this to the algebra of the product group of n cyclic groups of order 2 (cf.~§ 4, no. 1, Example 2).

\begin{enumerate}
\def\labelenumi{\arabic{enumi}.}
\setcounter{enumi}{7}
\item
  The quaternionic group \(\mathfrak{Q}\) (I, § 6, Exercise 4) is isomorphic to the multiplicative group of the eight quaternions \(\pm 1\) , \(\pm i\) , \(\pm j\) , \(\pm k\) in the quaternion algebra of type \((-1, -1)\) over a field in which \(2 \neq 0\) . Show that the algebra of the group \(\mathfrak{Q}\) over a field K in which \(2 \neq 0\) is isomorphic to the product of four fields isomorphic to K and the quaternion algebra of type \((+1, -1)\) over K (if \(c\) is the element of \(\mathfrak{Q}\) corresponding to the quaternion \(-1\) , the elements of \(\mathfrak{Q}\) can be written as \(e\) , \(i, j, k, c, ci, cj, ck\) ; consider the basis of E consisting of the elements \(\frac{1}{2}(e + c)\) , \(\frac{1}{2}(e - c)\) , \(\frac{1}{2}(e + c)i\) , \(\frac{1}{2}(e - c)i\) , \(\frac{1}{2}(e + c)j\) , \(\frac{1}{2}(e - c)j\) , \(\frac{1}{2}(e + c)k\) , \(\frac{1}{2}(e - c)k\) ).
\item
  \begin{enumerate}
  \def\labelenumii{(\alph{enumii})}
  \tightlist
  \item
    Show that the algebra E of the dihedral group \(\mathbf{D}_4\) of order 8 (I, §6, Exercise 4) over a field K in which \(2 \neq 0\) is isomorphic to the product of four fields isomorphic to K and the algebra of matrices of order 2 over K. (If \(a, b\) are the two generators of \(\mathbf{D}_4\) introduced in I, §6, Exercise 4, the elements of \(\mathbf{D}_4\) can be written as \(a^i b^j\) with \(0 \leqslant i \leqslant 3\) , \(0 \leqslant j \leqslant 1\) ; consider the basis of E consisting of the four elements \(\frac{1}{2}(e + a^2)\) , \(\frac{1}{2}(e - a^2)\) , \(\frac{1}{2}(a + a^3)\) , \(\frac{1}{2}(a - a^3)\) and these four elements multiplied on the right by \(b\) ; use Exercise 4.) Deduce that, if \(-1\) is a square in K, the algebras over K of the non-isomorphic groups \(\mathfrak{Q}\) and \(\mathbf{D}_4\) are isomorphic.
  \end{enumerate}
\end{enumerate}

\begin{enumerate}
\def\labelenumi{(\alph{enumi})}
\setcounter{enumi}{1}
\tightlist
\item
  Show similarly that the algebra F of the dihedral group \(D_{3}\) of order 6 over a field K in which \(2 \neq 0\) and \(3 \neq 0\) is isomorphic to the product of two fields isomorphic to K and the matrix algebra of order 2 over K (consider here the basis of F consisting of the elements \(e + a + a^{2}\) , \(a + a^{2} - 2e\) , \(a - a^{2}\) and these three elements multiplied on the right by b).
\end{enumerate}

\begin{enumerate}
\def\labelenumi{\arabic{enumi}.}
\setcounter{enumi}{9}
\item
  Let G be a group and H a normal subgroup of G. Show that, if a is the two-sided ideal of the group algebra \(\mathbf{A}^{(\mathbf{G})}\) , generated by the elements ts - s, where t runs through H and s runs through G, the group algebra \(\mathbf{A}^{(\mathbf{G}/\mathbf{H})}\) is isomorphic to the quotient algebra \(\mathbf{A}^{(\mathbf{G})}/\mathbf{a}\) .
\item
  Let A be a commutative ring and S a commutative monoid with identity element e. Suppose that an external law \((s, x) \mapsto x^{s}\) is given on A with S as domain of operators, such that, for all \(s \in S\) , the mapping \(x \mapsto x^{s}\) is an automorphism of the ring A and \((x^{s})^{t} = x^{st}\) for all \(x \in A\) and s and t in S (which implies that e is the identity operator for the external law considered). Then a multiplicative internal law of composition is defined on the A-module \(\mathrm{A}^{(\mathrm{S})}\) , whose canonical basis is denoted by \((b_{\mathrm{s}})\) , by the relation
\end{enumerate}

\[
\left(\sum_ {s} x _ {s} b _ {s}\right) \left(\sum_ {s} y _ {s} b _ {s}\right) = \sum_ {s} \left(\sum_ {t u = s} a _ {t, u} x _ {t} y _ {u} ^ {s}\right) b _ {s}
\]

where the \(a_{s,t}\) belong to A.

Show that this law and addition on \(\mathbf{A}^{(\mathbf{S})}\) define a ring structure on this set, provided the \(a_{s,t}\) satisfy the conditions

\[
a _ {s, t} a _ {s t, u} = a _ {t u} ^ {s} a _ {s, t u}
\]

for all s, t, u in S. The element \(b_{e}\) is unit element of this ring if also

\[
a _ {e, s} = a _ {s, e} = 1
\]

for all \(s \in S\) ; in that case A can be identified with a subring of \(A^{(S)}\) and if C denotes the subring of A consisting of the elements invariant under all the automorphisms \(x \mapsto x^s\) , \(A^{(S)}\) is a C-algebra, called the crossed product of the ring A and the monoid S, relative to the factor system \((a_{s,t})\) . If the factor system \((a_{s,t})\) is replaced by \(\left(\frac{c_s c_t^s}{c_{st}} a_{s,t}\right)\) , where, for all \(s \in S\) , \(c_s\) is an invertible element of A and \(c_e = 1\) , the new crossed product obtained is isomorphic to the crossed product defined by the factor system \((a_{s,t})\) .

If in particular A is taken to be a quadratic algebra over a ring C and S a cyclic group of order 2, say \(\{e, s\}\) , such that \(x^s = \bar{x}\) for \(x \in \mathbf{A}\) , show that every crossed product of A and S is a quaternion algebra over C.

\begin{enumerate}
\def\labelenumi{\arabic{enumi}.}
\setcounter{enumi}{11}
\tightlist
\item
  Let I be a non-empty set. Show that the elements of the free magma M(I) can be identified with the sequences \((a_{i})_{1 \leqslant i \leqslant n}\) (n an arbitrary integer \(\geqslant 1\) ) where the \(a_{i}\) are either elements of I or a special symbol P (representing the ``open brackets''), where these sequences satisfy the following conditions:
\end{enumerate}

\begin{enumerate}
\def\labelenumi{(\roman{enumi})}
\item
  The length \(n\) of the sequence is equal to twice the number of symbols \(\mathbf{P}\) which appear in it, plus 1.
\item
  For all \(k \leqslant n - 1\) , the number of symbols \(\mathbf{P}\) which appear in the subsequence \((a_i)_{1 \leqslant i \leqslant k}\) is \(\geqslant k / 2\) .
\end{enumerate}

(Use Set Theory, I, Appendix to associate with each sequence \((a_i)\) an element of M(I); for example, to the sequence PPPxyzPxy corresponds the word (((xy)z)(xy)).)

¶ 13. Let A be a commutative ring and E an A-module. An A-module \(E_{n}\) is defined inductively by writing

\[
\mathbf {E} _ {1} = \mathbf {E}, \quad \mathbf {E} _ {n} = \bigoplus_ {p + q = n} \left(\mathbf {E} _ {p} \otimes_ {\mathbf {A}} \mathbf {E} _ {q}\right) \quad \text { for } n \geqslant 2.
\]

We write \(\mathbf{ME} = \bigoplus_{n=1}^{\infty} \mathbf{E}_n\) and an A-algebra structure is defined on ME by setting, for \(x_p \in \mathbf{E}_p\) and \(x_q \in \mathbf{E}_q\) , \(x_p x_q = x_p \otimes x_q \in \mathbf{E}_{p+q}\) .

\begin{enumerate}
\def\labelenumi{(\alph{enumi})}
\item
  Show that for every A-algebra B and every A-linear mapping \(f: \mathbf{E} \to \mathbf{B}\) , there exists one and only one A-homomorphism of algebras \(\mathbf{ME} \to \mathbf{B}\) extending \(f\) . ME is called the free algebra of the A-module E.
\item
  An integer \(a(n)\) is defined by induction on \(n\) by the formulae
\end{enumerate}

\[
a (1), \quad a (n) = \sum_ {p + q = n} a (p) a (q) \quad \text { if } n \geqslant 2.
\]

Show that \(E_{n}\) is isomorphic to the direct sum of \(a(n)\) modules isomorphic to \(E^{\otimes n}\) .

\begin{enumerate}
\def\labelenumi{(\alph{enumi})}
\setcounter{enumi}{2}
\tightlist
\item
  Let \(f(\mathbf{T})\) be the formal power series \(\sum_{n=1}^{\infty} a(n)\mathbf{T}^n\) . Show that
\end{enumerate}

\[
f (\mathbf {T}) = \mathbf {T} + (f (\mathbf {T})) ^ {2}.
\]

*Deduce the formulae

\[
f (\mathrm{T}) = \frac {1 - \sqrt {1 - 4 \mathrm{T}}}{2} \quad \text { and } \quad a (n) = 2 ^ {n - 1} \frac {1 . 3 . 5 \dots (2 n - 1)}{n !}. *
\]

\begin{enumerate}
\def\labelenumi{(\alph{enumi})}
\setcounter{enumi}{3}
\tightlist
\item
  If I is any set and \(E = A^{(I)}\) , show that \(\operatorname{Lib}_{\mathbf{A}}(\mathbf{I})\) is identified with ME. There are canonical homomorphisms \(\mathbf{M}(\mathbf{I}) \xrightarrow{w} \mathbf{N}^{(\mathbf{I})} \xrightarrow{l} \mathbf{N}\) whose composition is length in \(\mathbf{M}(\mathbf{I})\) ; \(\omega(m)\) is called the multidegree of \(m \in \mathbf{M}(\mathbf{I})\) ; the mapping \(\omega\) defines on \(\operatorname{Lib}_{\mathbf{A}}(\mathbf{I})\) a graduation of type \(N^{(I)}\) . Show that the set \((\operatorname{Lib}_{\mathbf{A}}(\mathbf{I}))_{\alpha}\) of homogeneous elements of multidegree \(\alpha \in N^{(I)}\) under this graduation is a free A-module of rank equal to
\end{enumerate}

\[
a (\alpha) = a (n) \frac {n !}{\alpha !} = 2 ^ {n - 1} \frac {1 . 3 . 5 \dots (2 n - 1)}{\alpha !}
\]

where \(\alpha! = \prod_{i \in I} (\alpha(i))!\) and \(n = \sum_{i \in I} \alpha(i) = l(\alpha)\) is the length of \(\alpha\) .

\begin{enumerate}
\def\labelenumi{\arabic{enumi}.}
\setcounter{enumi}{13}
\tightlist
\item
  \begin{enumerate}
  \def\labelenumii{(\alph{enumii})}
  \tightlist
  \item
    Let M be a monoid whose law of composition is denoted by \(\top\) ; suppose further that M is totally ordered by an order relation \(x \leqslant y\) such that the relations \(x < y\) and \(x' \leqslant y'\) , or \(x \leqslant y\) and \(x' < y'\) , imply \(x \top x' < y \top y'\) . Show that if A is an integral domain, the algebra over A of the monoid M has no divisors of 0.
  \end{enumerate}
\end{enumerate}

\begin{enumerate}
\def\labelenumi{(\alph{enumi})}
\setcounter{enumi}{1}
\item
  Deduce that if \(\mathbf{A}\) is an integral domain, every polynomial algebra \(\mathbf{A}[(\mathbf{X}_i)_{i\in I}]\) is an integral domain.
\item
  If A is an integral domain and f and g two polynomials in \(\mathrm{A}[(\mathrm{X}_{i})_{i\in\mathbf{I}}]\) such that fg is homogeneous and \(\neq0\) , show that f and g are homogeneous. In particular, the invertible elements of the ring \(\mathrm{A}[(\mathrm{X}_{i})_{i\in\mathbf{I}}]\) are the invertible elements of the ring A.
\end{enumerate}

\begin{enumerate}
\def\labelenumi{\arabic{enumi}.}
\setcounter{enumi}{14}
\item
  Let A be a commutative ring and n an integer \(\geqslant2\) such that for all \(a\in A\) the equation nx=a has a solution \(x\in A\) . Let m be an integer \(\geqslant1\) and u a polynomial in A{[}X{]} of the form \(X^{mn}+f(X)\) where f is of degree \(\leqslant mn-1\) ; show that there exists a polynomial \(v\in A[X]\) of the form \(X^{m}+w\) , where w is of degree \(\leqslant m-1\) , such that \(u-v^{n}\) is zero or of degree \(<m(n-1)\) .
\item
  Let A be a commutative ring and \(u = \sum_{k=0}^{m} a_k X^k\) a divisor of 0 in the ring A{[}X{]}. Show that if \(v = \sum_{k=0}^{n} b_k X^k\) is an element \(\neq 0\) of A{[}X{]} of degree \(n\) such that \(uv = 0\) , there exists a polynomial \(w \neq 0\) of degree \(\leqslant n - 1\) such that \(uw = 0\) . (Reduce it to the case where \(b_0 \neq 0\) ; if \(a_k v = 0\) for \(0 \leqslant k \leqslant m - 1\) , show that we can take \(w = b_0\) ; if \(a_k v = 0\) for
\end{enumerate}

\[
0 \leqslant k <   p \leqslant m - 1
\]

and \(a_{p}v \neq 0\) , show that \(a_{p}b_{0} = 0\) and therefore we can take \(w = \sum_{k=0}^{n-1} a_{p}b_{k+1}\mathbf{X}^{k}\) . Deduce that there exists an element \(c \neq 0\) of \(\mathbf{A}\) such that \(cu = 0\) .

\begin{enumerate}
\def\labelenumi{\arabic{enumi}.}
\setcounter{enumi}{16}
\item
  Let A be a commutative ring. Show that, in the algebra A{[}X{]} of polynomials in one indeterminate over A, the mapping \((u,v)\mapsto u(v)\) is a law of composition which is associative and left distributive with respect to addition and multiplication in A{[}X{]}. If A is an integral domain, show that the relation \(u(v)=0\) implies that u=0 or that v is constant and that, if \(u\neq0\) and the degree of v is \textgreater0, the degree of \(u(v)\) is equal to the product of the degrees of u and v.
\item
  Let K be a commutative field and K{[}X{]} the ring of polynomials in one indeterminate over K. Show that every automorphism s of the ring K{[}X{]} leaves K invariant (Exercise 14(c)) and induces on K an automorphism of this field; show further that necessarily \(s(\mathbf{X}) = a\mathbf{X} + b\) , with \(a \neq 0\) in K and \(b \in \mathbf{K}\) (cf.~Exercise 17). Conversely, show that being given an automorphism \(\sigma\) of \(\mathbf{K}\) and two elements \(a, b\) of \(\mathbf{K}\) such that \(a \neq 0\) defines one and only one automorphism \(s\) of \(\mathbf{K}[\mathbf{X}]\) such that \(s|_{\mathbf{K}} = \sigma\) and \(s(\mathbf{X}) = a\mathbf{X} + b\) . If \(G\) is the group of all automorphisms of the ring \(\mathbf{K}[\mathbf{X}]\) and \(N\) the subgroup of \(G\) consisting of the K-algebra structure of \(\mathbf{K}[\mathbf{X}]\) , show that \(N\) is a normal subgroup of \(G\) and that \(G / N\) is isomorphic to the automorphism group of the field \(K\) ; the group \(N\) is isomorphic to the group defined on the set \(K^* \times K\) by the law of composition \((\lambda, \mu)(\lambda', \mu') = (\lambda\lambda', \lambda' \mu + \mu')\) .
\item
  Prove the polynomial identity in 8 indeterminates
\end{enumerate}

\[
\begin{array}{r l} & (\mathrm{X} _ {1} ^ {2} + \mathrm{X} _ {2} ^ {2} + \mathrm{X} _ {3} ^ {2} + \mathrm{X} _ {4} ^ {2}) (\mathrm{Y} _ {1} ^ {2} + \mathrm{Y} _ {2} ^ {2} + \mathrm{Y} _ {3} ^ {2} + \mathrm{Y} _ {4} ^ {2}) \\ & \qquad = (\mathrm{X} _ {1} \mathrm{Y} _ {1} - \mathrm{X} _ {2} \mathrm{Y} _ {2} - \mathrm{X} _ {3} \mathrm{Y} _ {3} - \mathrm{X} _ {4} \mathrm{Y} _ {4}) ^ {2} \\ & \qquad + (\mathrm{X} _ {1} \mathrm{Y} _ {2} + \mathrm{X} _ {2} \mathrm{Y} _ {1} + \mathrm{X} _ {3} \mathrm{Y} _ {4} - \mathrm{X} _ {4} \mathrm{Y} _ {3}) ^ {2} \\ & \qquad + (\mathrm{X} _ {1} \mathrm{Y} _ {3} + \mathrm{X} _ {3} \mathrm{Y} _ {1} + \mathrm{X} _ {4} \mathrm{Y} _ {2} - \mathrm{X} _ {2} \mathrm{Y} _ {4}) ^ {2} \\ & \qquad + (\mathrm{X} _ {1} \mathrm{Y} _ {4} + \mathrm{X} _ {4} \mathrm{Y} _ {1} + \mathrm{X} _ {2} \mathrm{Y} _ {3} - \mathrm{X} _ {3} \mathrm{Y} _ {2}) ^ {2}. \end{array}
\]

(Apply formula (32) of no. 5 to a suitable quaternion algebra.)

\begin{enumerate}
\def\labelenumi{\arabic{enumi}.}
\setcounter{enumi}{19}
\tightlist
\item
  Show that there exists no polynomial identity in 6 indeterminates
\end{enumerate}

\[
(\mathrm{X} _ {1} ^ {2} + \mathrm{X} _ {2} ^ {2} + \mathrm{X} _ {3} ^ {2}) (\mathrm{Y} _ {2} ^ {2} + \mathrm{Y} _ {2} ^ {2} + \mathrm{Y} _ {3} ^ {2}) = \mathrm{P} ^ {2} + \mathrm{Q} ^ {2} + \mathrm{R} ^ {2}
\]

where P, Q, R are three polynomials with coefficients in Z with respect to the indeterminates \(X_{i}\) and \(Y_{i}\) . (Observe that 15 = 3.5 cannot be written in the form \(p^{2} + q^{2} + r^{2}\) , where p, q, r are integers.)

\begin{enumerate}
\def\labelenumi{\arabic{enumi}.}
\setcounter{enumi}{20}
\tightlist
\item
  Let S be a multiplicative monoid with identity element e, satisfying condition (D) of no. 10 and further satisfying the two following conditions: (1) the relation st = e implies s = t = e: (2) for all \(s \in S\) , the number n of terms in a finite sequence \((t_{i})_{1 \leq i \leq n}\) of elements distinct from e and such that \(t_{1}t_{2} \ldots t_{n} = s\) is less than a finite number \(v(s)\) depending only on s.
\end{enumerate}

Show that, for an element \(x = \sum_{s \in S} \alpha_s s\) of the total algebra of \(S\) over a ring \(A\) to be invertible, it is necessary and sufficient that \(\alpha_e\) be invertible in \(A\) (reduce it to the case where \(\alpha_e = 1\) and use the identity

\[
e - z ^ {n + 1} = (e - z) (e + z + \dots + z ^ {n})).
\]

\begin{enumerate}
\def\labelenumi{\arabic{enumi}.}
\setcounter{enumi}{21}
\item
  If K is a commutative field, show that in the ring of formal power series \(K[[X_{1}, X_{2}, \ldots, X_{p}]]\) there is only a single maximal ideal, equal to the set of non-invertible elements. On the other hand, show that in the polynomial ring \(K[X_{1}, \ldots, X_{p}]\) for \(p \geqslant 1\) there always exist several distinct maximal ideals.
\item
  Let K be a commutative field; show that there exists no formal power series \(u(\mathbf{X}, \mathbf{Y}) \in \mathbf{K}[[\mathbf{X}, \mathbf{Y}]]\) such that, for an integer \(m > 0\) ,
\end{enumerate}

\[
(\mathrm{X} + \mathrm{Y}) u (\mathrm{X}, \mathrm{Y}) = \mathrm{X} ^ {m} \mathrm{Y} ^ {m}.
\]

\begin{enumerate}
\def\labelenumi{\arabic{enumi}.}
\setcounter{enumi}{23}
\item
  Let K be a commutative field and let k be an integer such that \(k \neq 0\) in K. Show that for every formal power series u of K{[}{[}X{]}{]} with constant term equal to 1, there exists a formal power series \(v \in K[[X]]\) such that \(v^{k} = u\) .
\item
  Let A be a ring, commutative or otherwise, and let \(\sigma\) be an endomorphism of A. On the additive product group \(\mathbf{E} = \mathbf{A}^{\mathbf{N}}\) an internal law of composition is defined by writing \((\alpha_{n})(\beta_{n}) = (\gamma_{n})\) , where \(\gamma_{n} = \sum_{p + q = n}\alpha_{p}\sigma^{p}(\beta_{q})\) (with the convention \(\sigma^0 (\xi) = \xi\) ).
\end{enumerate}

\begin{enumerate}
\def\labelenumi{(\alph{enumi})}
\tightlist
\item
  Show that this law defines, with addition on E, a ring structure on E, admitting as unit element e the sequence \((\alpha_{n})\) where \(\alpha_{0}=1\) and \(\alpha_{n}=0\) for \(n\geqslant1\) . The mapping which associates with every \(\xi\in A\) the element \((\alpha_{n})\in E\) , such that \(\alpha_{0}=\xi,\alpha_{n}=0\) for \(n\geqslant1\) , is an isomorphism of A onto a
\end{enumerate}

subring of E, with which it is identified. Then we write \(\sum_{n=0}^{\infty} \alpha_n X^n\) instead of

\[
u = \sum_ {n = 0} ^ {\infty} \alpha_ {n} X ^ {n} \neq 0.
\]

\((\alpha_{n})\) and \(\mathbf{X}^{p}\beta = \sigma^{p}(\beta)\mathbf{X}^{p}\) for all \(\beta \in \mathbf{A}\) ; if \(u = \sum_{n=0}^{\infty} \alpha_{n}\mathbf{X}^{n} \neq 0\) , the smallest integer \(n\) such that \(\alpha_{n} \neq 0\) is called the order of \(u\) and denoted by \(\omega(u)\) .

\begin{enumerate}
\def\labelenumi{(\alph{enumi})}
\setcounter{enumi}{1}
\tightlist
\item
  If A is a ring with no divisor of 0 and \(\sigma\) is an isomorphism of A onto a subring of A, show that E is a ring with no divisor of 0 and that
\end{enumerate}

\[
\omega (u v) = \omega (u) + \omega (v)
\]

for \(u \neq 0\) and \(v \neq 0\) .

\begin{enumerate}
\def\labelenumi{(\alph{enumi})}
\setcounter{enumi}{2}
\item
  For \(u = \sum_{n=0}^{\infty} \alpha_n X^n\) to be invertible, it is necessary and sufficient that \(\alpha_0\) be invertible in A.
\item
  Suppose that A is a field and \(\sigma\) an automorphism of A. Show that E admits a field of left fractions (I, §9, Exercise 15) and that in this field F every element \(\neq0\) can be written uniquely in the form \(uX^{-h}\) , where u is an element of E of order 0.
\end{enumerate}

¶ *26. (a) Let A and B be two well-ordered subsets of R (which are necessarily countable; cf.~General Topology, IV, § 2, Exercise 1). Show that A + B is well-ordered and that, for all \(c \in \mathbf{A} + \mathbf{B}\) , there exists only a finite number of ordered pairs \((a, b)\) such that \(a \in \mathbf{A}\) , \(b \in \mathbf{B}\) and \(a + b = c\) (to show that a non-empty subset of A + B has a least element, consider its greatest lower bound in R).

\begin{enumerate}
\def\labelenumi{(\alph{enumi})}
\setcounter{enumi}{1}
\tightlist
\item
  Let K be a commutative field. In the vector space \(K^{R}\) , consider the vector subspace E consisting of the elements \((\alpha_{x})\) such that the set (depending on \((\alpha_{x})\) of \(x \in \mathbf{R}\) such that \(\alpha_{x} \neq 0\) is well-ordered. For any two elements \((\alpha_{x})\) , \((\beta_{x})\) of E, we write \((\alpha_{x})(\beta_{x}) = (\gamma_{x})\) , where \(\gamma_{x} = \sum_{y+z=x} \alpha_{y} \beta_{z}\) (a sum which is meaningful by virtue of (a)). Show that this law of composition defines, with addition, a field structure on E; the elements of E are also denoted by \(\sum_{t \in \mathbb{R}} \alpha_{t} X^{t}\) and called formal power series with well-ordered real exponents, with coefficients in K.*
\end{enumerate}

